
%(BEGIN_QUESTION)
% Copyright 2014, Tony R. Kuphaldt, released under the Creative Commons Attribution License (v 1.0)
% This means you may do almost anything with this work of mine, so long as you give me proper credit

Les øvelse 6 (``Creating a PID Control Loop (FIC-101) from Scratch'') i kapittel 4 (``Creating and Downloading the Control Strategy'') i manualen ``Getting Started with your DeltaV Digital Automation System'' (dokument D800002X122, mars 2006), og svar på følgende spørsmål:

\vskip 10pt

I denne øvelsen vises hvordan brukeren kan begynne å konfigurere en reguleringsmodul uten å starte helt fra bunnen, slik som i øvelse 5. Hvordan gjøres dette?

\vskip 10pt

Vanligvis kobles {\tt PID}-funksjonsblokker til analog inngangsblokk ({\tt AI}) og analog utgangsblokk ({\tt AO}) for å lage en komplett sløyfemodul. I denne øvelsen står {\tt PID}-blokken alene. I stedet for {\tt AI}- og {\tt AO}-blokker som ruter signaler til/fra fysisk I/O, hvordan ``vet'' denne {\tt PID}-blokken hvor den skal hente PV-inngang og hvor MV-utgang skal sendes?



\vskip 20pt \vbox{\hrule \hbox{\strut \vrule{} {\bf Forslag til sokratisk diskusjon} \vrule} \hrule}

\begin{itemize}
\item{} Gå til en DeltaV-arbeidsstasjon og åpne en modul i Control Studio. Finn noen av {\tt PID}-parametrene og valgene omtalt i øvelsen. {\it Ikke ``download'' eller ``save'' noe som endrer DCS-konfigurasjonen -- bare utforsk og observer!}
\item{} Noter hvor regulatorens virkningsretning (``direct'' eller ``reverse'') velges i {\tt PID}-blokken. Hvordan avgjør man riktig virkningsretning for en konkret prosess?
\item{} Rett etter innstillingen av regulatorretning beskriver øvelsen en lignende parameter i IO\_OPTS kalt {\it Increase to Close}. Den brukes når reguleringsventilen er luft-for-å-lukke (feiler åpen), slik at utgangsstolpen i operatørbildet samsvarer med ventilstilling (0\% = stengt ventil, 100\% = helt åpen). Forklar hvorfor denne parameteren er viktig å sette riktig i slike prosesser. En god forklaring er å beskrive hva som skjer hvis parameteren settes feil.
\item{} Sammenlign innstilling av prosessvariabelens ``engineering units'' med MINSCALE og MAXSCALE i {\tt AI}-blokken i Siemens 353. Hvordan er disse oppgavene like, og hvordan er de forskjellige?
\end{itemize}


\underbar{file i00814}
%(END_QUESTION)





%(BEGIN_ANSWER)


%(END_ANSWER)





%(BEGIN_NOTES)

På side 4-27 instrueres brukeren til å kopiere en mal fra DeltaV-biblioteket, gi den kopierte modulen nytt navn, og redigere videre derfra.

\vskip 10pt

På side 4-28 instrueres leseren til å sette parameterne {\tt IO\_IN} og {\tt IO\_OUT} i {\tt PID}-funksjonsblokken, som angir hvor blokken henter inngangssignal (PV) fra og hvor utgangssignal (MV) sendes.

%INDEX% Reading assignment: Emerson DeltaV "Getting Started" manual (Chapter 4, Exercise 6)

%(END_NOTES)
